\documentclass{training}
\schoollogo{../../../resources/latex/crack.pdf}
\usepackage[utf8]{inputenc}
\usepackage[T1]{fontenc}
\usepackage{lmodern}
\setlang{french}
\usepackage{hyperref}
\usepackage{graphicx}
\usepackage{tabularx}
\usepackage{fvextra}
\fvset{breaklines=true,breakanywhere=true,breaksymbol={},breaksymbolleft={},breaksymbolright={},breakanywheresymbolpre={},breakanywheresymbolpost={}}
\makeatletter
\g@addto@macro\UrlBreaks{%
  \do\a\do\b\do\c\do\d\do\e\do\f\do\g\do\h\do\i\do\j\do\k\do\l\do\m%
  \do\n\do\o\do\p\do\q\do\r\do\s\do\t\do\u\do\v\do\w\do\x\do\y\do\z%
  \do\A\do\B\do\C\do\D\do\E\do\F\do\G\do\H\do\I\do\J\do\K\do\L\do\M%
  \do\N\do\O\do\P\do\Q\do\R\do\S\do\T\do\U\do\V\do\W\do\X\do\Y\do\Z%
  \do\0\do\1\do\2\do\3\do\4\do\5\do\6\do\7\do\8\do\9%
  \do\<\do\>\do\(\do\)\do\'\do\/}
\schoollogo{../../../resources/latex/crack.pdf}
\makeatother


\begin{document}

\title{nf\_tables : chasse aux variantes --- A-EPI-002}
\subtitle{Générations, catchall, GC}
\member{Ted ``Tedsig42'' Kouhouenou}{ted.sig42@tutamail.com}
\summary{Module avancé : modéliser les transactions nf\_tables, traquer les variantes de quatre CVE, écrire des règles Semgrep/Weggli pour les motifs de génération et trier les résultats sur les sources du noyau.}
\version{1.0}
\prerequisite{A-EPI-001, lecture de C noyau, notions de RCU et de GC}
\maketitle

\newpage
\tableofcontents
\newpage

\chapter{Instructions}
\begin{itemize}
  \item Deux à trois jours.
  \item Travail sur une copie locale du noyau, jamais sur une cible de
        production.
  \item Les règles doivent être testées sur un petit corpus, pas lancées
        sur tout le noyau sans filtre.
  \item Aucun exploit, aucun module noyau chargeable sur une machine
        partagée.
\end{itemize}

\subsection*{Format du rendu --- en-tête obligatoire}

Tout fichier rendu dans ce module doit commencer par l'en-tête
ci-dessous, \textbf{adapté à la syntaxe de commentaire du langage}
utilisé.

\subsubsection*{Modèle générique}

\begin{minted}{c}
/**
 * Module: <code module>
 * Exercise: <numéro exo>
 * File: <nom du fichier>
 * Author: <votre login JeuneCrack>
 *
 * Description: <une phrase, pas de roman>
 *
 * Return: <voir ci-dessous>
 */
\end{minted}

\subsubsection*{Login}

Votre identifiant \textbf{JeuneCrack} (ex: \texttt{tedsig42}).
Pour un rendu en binôme : \texttt{login1/login2}.
Ajoutez un handle Matrix après le login si vous le souhaitez
(ex: \texttt{tedsig42 @tedsig42}).

\subsubsection*{Return --- selon le type de fichier}

\begin{tabularx}{\linewidth}{lX}
  \toprule
  \textbf{Type de fichier} & \textbf{Valeur de Return} \\
  \midrule
  C (\texttt{.c}) &
  \texttt{return (0);} (succès),
  \texttt{return (-1);} (erreur),
  \texttt{return (valeur);} (résultat calculé) \\

  Python (\texttt{.py}) &
  \texttt{0} (succès),
  \texttt{1} (erreur),
  \texttt{N/A} (analyse, script sans exit) \\

  Shell (\texttt{.sh}) &
  \texttt{exit 0} (succès),
  \texttt{exit 1} (erreur) \\

  Markdown (\texttt{.md}) &
  \texttt{N/A} --- les fichiers d'analyse n'ont pas de valeur de retour \\

  Exploit C/Python &
  \texttt{shell} (shell obtenu),
  \texttt{flag: <flag>} (flag capturé),
  \texttt{proof: <hash>} (preuve de concept) \\

  Binaire (\texttt{.bin}) &
  \texttt{N/A} --- pas d'en-tête pour les fichiers binaires \\
  \bottomrule
\end{tabularx}

\subsubsection*{Exemples par langage}

\textbf{C (\texttt{.c})} :
\begin{minted}{c}
/**
 * Module: F-EXP-001
 * Exercise: ex02
 * File: exploit.c
 * Author: tedsig42
 *
 * Description: Buffer overflow with ret2libc
 *
 * Return: 0
 */
\end{minted}

\textbf{Python (\texttt{.py})} :
\begin{minted}{python}
# Module: F-WEB-001
# Exercise: ex05
# File: ssrf-exploit.py
# Author: tedsig42
#
# Description: Blind SSRF bypass via redirect
#
# Return: 0
\end{minted}

\textbf{Shell (\texttt{.sh})} :
\begin{minted}{bash}
# Module: F-SUP-001
# Exercise: ex03
# File: dependency-confusion.sh
# Author: tedsig42
#
# Description: Automated dependency confusion test
#
# Return: exit 0
\end{minted}

\textbf{Markdown (\texttt{.md})} --- analyse, rapport :
\begin{Verbatim}[fontsize=\normalsize]
<!--
Module: M-SRC-001
Exercise: ex04
File: audit-report.md
Author: tedsig42

Description: Code audit of vulnerable-httpd

Return: N/A
-->
\end{Verbatim}

\chapter{Instructions Communes}

\section{Structure du D\'ep\^ot}

Votre d\'ep\^ot Git doit respecter la convention de nommage suivante pour
tous les modules soumis. L'outil d'\'evaluation (\texttt{pendragon})
se base sur cette structure pour localiser et tester votre travail.

\begin{console}
<login>/
+-- <MODULE-ID>-<slug>/         # ex. M-BLS-001-business-logic/
|   +-- <turnindir>/            # ex. ex00-pricing-manipulation/
|   |   +-- <file1>             # ex. poc-pricing.py
|   |   +-- <file2>             # ex. report.md
|   |   +-- ...
|   +-- <turnindir>/            # ex. ex01-race-conditions/
|   |   +-- ...
|   +-- ...
+-- ...
\end{console}

\begin{itemize}
  \item \textbf{R\'epertoire du module} : le nom complet du r\'epertoire tel
        qu'affich\'e dans votre curriculum (ex. \texttt{M-BLS-001-business-logic}).
  \item \textbf{R\'epertoire de rendu} : la valeur de \texttt{\textbackslash turnindir}
        fournie dans chaque exercice (ex. \texttt{ex00-pricing-manipulation}).
  \item \textbf{Fichiers de l'exercice} : exactement les fichiers list\'es dans
        \texttt{\textbackslash exfiles}. Ni fichiers suppl\'ementaires, ni fichiers manquants.
\end{itemize}

\section{Qu'est-ce qu'une Preuve de Concept Valide}

Une PoC doit \^etre :

\begin{enumerate}
  \item \textbf{Reproductible}. N'importe qui avec le m\^eme environnement peut
        l'ex\'ecuter et observer la vuln\'erabilit\'e. Les chemins cod\'es en dur,
        les adresses IP absolues et les jetons expir\'es ne sont pas acceptables.
  \item \textbf{Document\'ee}. Chaque \'etape est expliqu\'ee. L'\'evaluteur ne
        devrait pas avoir \`a deviner ce que fait le script.
  \item \textbf{Cibl\'ee}. L'exercice sp\'ecifie une cible autoris\'ee
        (lab, maquette, binaire local ou plateforme d\'esign\'ee). Ne testez
        jamais les PoC sur des syst\`emes de production qui ne vous appartiennent pas.
  \item \textbf{Minimale}. La PoC d\'emontre \emph{uniquement} la vuln\'erabilit\'e
        en question. Les charges utiles non li\'ees et le bruit sont p\'enalis\'es.
\end{enumerate}

\begin{warning}
  Une PoC qui modifie silencieusement l'\'etat du syst\`eme, supprime des
  fichiers ou persiste des changements en dehors du p\'erim\`etre de l'exercice
  sera rejet\'ee. Si votre exploit n\'ecessite un nettoyage, documentez les
  \'etapes de nettoyage dans votre rapport.
\end{warning}

\section{Processus de Soumission}

\begin{enumerate}
  \item R\'ealisez les livrables de l'exercice sur votre machine locale.
  \item Committez et poussez vers votre d\'ep\^ot Git assign\'e.
  \item Le d\'ep\^ot est gel\'e apr\`es la date limite.
  \item Les pairs clonent votre d\'ep\^ot et \'evaluent en utilisant :
        \begin{itemize}
          \item \texttt{pendragon} pour l'ex\'ecution automatis\'ee des tests
                (l\`a o\`u le module fournit des configurations mouli).
          \item La grille manuelle pour la qualit\'e du rapport, le style de
                code et la m\'ethodologie.
        \end{itemize}
\end{enumerate}

\hint{Ex\'ecutez toujours \texttt{pendragon run <MODULE-ID> <ex\_id>
<path/to/student/repo>} avant votre commit final pour v\'erifier
que votre soumission passe les v\'erifications automatis\'ees.}

\subsection*{Test avec pendragon}

Si le module inclut une configuration mouli (fichiers JSON dans
\texttt{resources/mouli/}), vous pouvez tester votre soumission localement :

\begin{console}
pendragon run M-BLS-001 ex00 ~/student-repos/tedsig42
\end{console}

\texttt{pendragon} r\'esout le chemin de l'exercice automatiquement en
cherchant \texttt{<mod-fullname>/<ex-dirname>/} dans le r\'epertoire
\'etudiant donn\'e.

\subsection*{V\'erification de style avec caliban}

Pour les modules qui n\'ecessitent du code C, \texttt{caliban} v\'erifie
la conformit\'e du style :

\begin{console}
caliban check ~/student-repos/tedsig42/M-BLS-001-business-logic
\end{console}

\texttt{caliban} applique la Norme 42 (taille des fonctions, nommage
des variables, constructions interdites) adapt\'ee aux contextes de
recherche de vuln\'erabilit\'es. Ex\'ecutez-le avant la soumission ---
un \'echec aux v\'erifications de style peut invalider l'exercice.

\section{R\`egle des Fichiers Suppl\'ementaires Interdits}

\begin{itemize}
  \item Votre d\'ep\^ot doit contenir \textbf{uniquement} les fichiers
        list\'es dans \texttt{\textbackslash exfiles} pour chaque exercice.
  \item Fichiers interdits : \texttt{.DS\_Store}, \texttt{\_\_pycache\_\_},
        \texttt{.venv/}, \texttt{node\_modules/}, \texttt{.git/},
        binaires compil\'es (sauf explicitement list\'es dans \texttt{exfiles}),
        fichiers d'\'echange d'\'editeur, fichiers journaux et tout fichier
        non demand\'e.
  \item Utilisez \texttt{.gitignore} pour exclure automatiquement les
        fichiers ind\'esirables.
\end{itemize}

\section{Format du Rapport}

Lorsqu'un rapport (\texttt{report.md}) est requis, il doit inclure :

\begin{enumerate}
  \item \textbf{Classe de vuln\'erabilit\'e} (identifiant CWE si applicable).
  \item \textbf{Étapes de reproduction} --- num\'erot\'ees, pr\^etes
        \`a \^etre copi\'ees-coll\'ees.
  \item \textbf{Charges utiles} --- chaînes exactes, requ\^etes ou
        entr\'ees utilis\'ees.
  \item \textbf{Impact} --- ce qu'un attaquant peut accomplir.
  \item \textbf{Mitigation} --- comment le d\'eveloppeur devrait corriger.
  \item \textbf{Sources} --- toutes les ressources consult\'ees (voir
        Section~\ref{sec:cite-sources}).
\end{enumerate}

\subsection*{En-t\^ete de Rendu}

Chaque fichier source doit inclure l'en-t\^ete de rendu d\'efini dans
\path{../../../resources/latex/deliverable-header.tex}, adapt\'e
\`a la syntaxe de commentaire du langage utilis\'e.

\section{Cibles Autoris\'ees Uniquement}

\begin{warning}
  Vous pouvez \textbf{uniquement} tester les cibles explicitement list\'ees
  dans l'exercice (labs, serveurs locaux, maquettes, binaires d\'esign\'es).
  Tester des syst\`emes de production r\'eels, des applications tierces
  ou toute cible qui ne vous appartient pas ou pour laquelle vous n'avez
  pas d'autorisation \'ecrite est interdit et peut entraîner une
  disqualification imm\'ediate.
\end{warning}

L'\'ethique du bug bounty s'applique tout au long de ce curriculum :
\begin{itemize}
  \item Aucun test sur des programmes en production sans un accord de
        p\'erim\`etre valide.
  \item Aucun d\'eni de service sauf si l'exercice l'exige explicitement.
  \item Aucune exfiltration de donn\'ees au-del\`a de ce qui est
        n\'ecessaire pour d\'emontrer la vuln\'erabilit\'e.
  \item Si vous acc\'edez accidentellement \`a des donn\'ees r\'eelles
        d'utilisateurs, arr\^etez imm\'ediatement et documentez ce qui
        s'est pass\'e.
\end{itemize}

\section{Citez Toutes Vos Sources}
\label{sec:cite-sources}

Toute ressource consult\'ee pendant votre travail doit \^etre cit\'ee.
Cela inclut :

\begin{itemize}
  \item Articles acad\'emiques et billets de blog.
  \item Rapports publics de bug bounty (HackerOne, Bugcrowd, etc.).
  \item Documentation officielle (RFC, pages de manuel, documentation
        fournisseur).
  \item \textbf{Outils d'IA} : nom du mod\`ele, r\'esum\'e de la
        question, date et m\'ethode de v\'erification (voir le chapitre
        Instructions IA pour le format de citation exact).
  \item Extraits de code adapt\'es de projets open-source (avec
        attribution de licence si requis).
\end{itemize}

\subsection*{Format de citation}

\begin{minted}{markdown}
# R\'ef\'erences

- [1] PortSwigger, "Business Logic Flaws",
  https://portswigger.net/web-security/logic-flaws
  (consult\'e le 2025-03-15).
- [2] Rapport HackerOne #123456, "IDOR on /api/users",
  anonymis\'e, consult\'e via divulgation publique.
- [3] G\'en\'er\'e avec l'assistance de Claude 3.5 Sonnet
  le 2025-03-14. Question : "Write a Python script that...".
  V\'erifi\'e par : revue manuelle de chaque requ\^ete HTTP.
\end{minted}

\subsection*{Attribution en ligne}

Si un exercice vous demande explicitement d'utiliser une technique ou
un outil sp\'ecifique, vous n'avez pas besoin de citer l'exercice
lui-m\^eme comme source. Citez tout ce qui se trouve en dehors de
l'\'enonc\'e de l'exercice.

\chapter{Instructions sur l'IA}

\section{Aperçu de la Politique}

Les outils de codage assistés par IA (ChatGPT, Claude, Copilot, Gemini,
Cursor, Codeium, et similaires) sont \textbf{autorisés} tout au long
de ce cursus. La recherche en sécurité est un domaine où la maîtrise
de l'IA devient essentielle --- apprendre à utiliser ces outils
efficacement fait partie de votre formation.

Cependant, l'objectif est la \textbf{compréhension}, pas l'automatisation.
Si vous ne pouvez pas expliquer chaque ligne de votre exploit sans
l'aide de l'IA, vous n'avez pas terminé l'exercice.

\begin{warning}
  Utiliser l'IA pour contourner l'objectif pédagogique d'un exercice
  --- par exemple, lui demander de rétro-ingéniérer un binaire pour
  ne pas avoir à apprendre le désassembleur, ou de générer une chaîne
  d'exploitation complète que vous ne comprenez pas --- va à l'encontre
  du but du cursus et sera traité comme une fraude académique.
\end{warning}

\section{Utilisations Autorisées}

\begin{itemize}
  \item \textbf{Explication de concepts}. Demander à l'IA de clarifier
        une classe de vulnérabilité, un sous-système du noyau, ou un
        protocole réseau.
  \item \textbf{Échafaudage de code}. Générer du code de départ pour
        un harnais, un fuzzer, ou un squelette de PoC --- puis le
        modifier et l'étendre vous-même.
  \item \textbf{Débogage}. Expliquer des messages d'erreur, la sortie
        d'un désassembleur, ou un comportement de crash inattendu.
  \item \textbf{Résumé de documentation}. Condenser une longue RFC ou
        un article technique en notes exploitables.
  \item \textbf{Traduction}. Traduire du contenu technique entre
        l'anglais et le français (les deux langues sont utilisées
        dans ce cursus).
  \item \textbf{Rédaction de rapports}. Demander de l'aide pour
        structurer un rapport de bug bounty --- mais remplir le
        contenu technique vous-même.
\end{itemize}

\section{Utilisations Interdites}

\begin{itemize}
  \item \textbf{Génération complète d'exploit}. Demander à l'IA
        d'écrire l'exploit complet d'un exercice, puis le soumettre
        sans le comprendre.
  \item \textbf{Contournement des contraintes de l'exercice}. Utiliser
        l'IA pour contourner une restriction que l'exercice impose
        explicitement (p.\ ex.\ «~n'utilisez pas de scanners
        automatisés~»).
  \item \textbf{Découverte automatisée de vulnérabilités}. Exécuter des
        outils de scan pilotés par IA qui remplacent l'analyse manuelle
        que vous êtes censé pratiquer.
  \item \textbf{Rédaction des grilles d'évaluation}. Vous n'êtes pas
        autorisé à utiliser l'IA pour générer des critères d'évaluation
        par les pairs ou des réponses.
  \item \textbf{Plagiat sans attribution}. Copier le résultat de l'IA
        sans le citer (voir Section~\ref{sec:ai-citation}).
\end{itemize}

\section{Responsabilité Personnelle}

Lorsque vous demandez à l'IA de générer du code d'exploit, vous
devenez responsable de chaque octet de ce résultat avant de
l'utiliser contre une cible.

\begin{enumerate}
  \item \textbf{Lisez le code généré} ligne par ligne.
  \item \textbf{Identifiez toutes les hypothèses} faites par l'IA
        (architecture, version du noyau, drapeaux du compilateur,
        topologie réseau).
  \item \textbf{Testez en isolation} avant d'exécuter contre une
        cible d'exercice.
  \item \textbf{Comprenez les modes d'échec}. Si l'exploit généré
        plante ou se bloque, vous devez diagnostiquer et corriger
        le problème vous-même.
\end{enumerate}

\hint{Si l'IA génère du code que vous ne comprenez pas, demandez-lui
de l'expliquer. Si vous ne comprenez toujours pas après
l'explication, le code n'est pas prêt à être utilisé. Passez à
une approche plus simple que vous maîtrisez entièrement.}

\section{Citation Obligatoire de l'IA}
\label{sec:ai-citation}

Chaque fichier contenant du code généré ou assisté par l'IA doit
inclure un commentaire de citation. Le format est~:

\begin{minted}{python}
# Generated with assistance from <tool> on <date>.
# Prompt: <one-line summary of the prompt>
# Verified by: <method used to verify the output>
\end{minted}

\subsection*{Exemples}

\textbf{Python}~:
\begin{minted}{python}
# Generated with assistance from Claude 3.5 Sonnet on 2025-03-14.
# Prompt: "Write a Python script that sends HTTP requests with
# modified price fields to /api/checkout/pay"
# Verified by: manual review of each request and response
\end{minted}

\textbf{C}~:
\begin{minted}{c}
/* Generated with assistance from GitHub Copilot on 2025-03-14.
 * Prompt: "Implement a stack-based buffer overflow with
 * ret2libc for x86_64 Linux"
 * Verified by: compiled with -fno-stack-protector and
 * tested against local binary with ASLR disabled */
\end{minted}

\textbf{Markdown (rapport)}~:
\begin{minted}{html}
<!--
Generated with assistance from ChatGPT-4 on 2025-03-14.
Prompt: "Structure a bug bounty report for IDOR with
sections: summary, steps, impact, mitigation"
Verified by: all technical content written manually,
AI only suggested structure
-->
\end{minted}

\subsection*{Interactions multiples avec l'IA}

Si vous avez utilisé plusieurs outils d'IA, ou itéré à travers
plusieurs invites, citez chaque interaction significative. Un
groupe d'invites connexes peut être résumé~:

\begin{minted}{html}
<!--
AI assistance used in this file:
- Claude 3.5 Sonnet (2025-03-14): harness skeleton generation
- ChatGPT-4 (2025-03-14): crash triage help
- Copilot (2025-03-15): inline autocomplete for loop logic
All generated code reviewed and modified manually.
-->
\end{minted}

\section{Évaluation par les Pairs et l'IA}

Lors de l'évaluation par les pairs, l'évaluateur peut~:

\begin{itemize}
  \item Vous demander d'\textbf{expliquer tout code généré par l'IA}
        ligne par ligne.
  \item Vous demander de \textbf{reproduire le raisonnement} sans
        l'aide de l'IA.
  \item Demander à voir votre \textbf{historique d'invites} pour un
        exercice donné.
  \item Déduire des points si la citation de l'IA est manquante ou
        si vous ne pouvez pas expliquer le code que l'IA a généré.
\end{itemize}

\begin{warning}
  Présenter le code généré par l'IA comme votre propre travail
  (défaut de citation) est du plagiat. Si l'évaluateur soupçonne
  du plagiat, il peut saisir l'équipe pédagogique pour examen.
\end{warning}

\section{Rappel Éthique}

\begin{itemize}
  \item Les modèles d'IA peuvent produire du code incorrect,
        dangereux, ou illégal. Vous êtes \textbf{responsable} de
        vérifier chaque résultat avant exécution.
  \item Les modèles d'IA entraînés sur la recherche publique en
        sécurité peuvent reproduire textuellement du code provenant
        de projets sous licence GPL. Vérifiez les licences avant
        d'intégrer ces éléments dans votre travail.
  \item Ne saisissez pas de code sensible ou propriétaire dans des
        outils d'IA hébergés dans le cloud. Utilisez des modèles
        locaux (Llama, Mistral, DeepSeek) pour les travaux impliquant
        des cibles non publiques.
  \item L'éthique de la recherche en sécurité que vous apprenez ici
        s'applique à l'utilisation de l'IA~: ne demandez pas à l'IA
        de générer des malwares, des rançongiciels, ou du code destiné
        à nuire à des systèmes en dehors du cadre de l'exercice.
\end{itemize}

\info{La maîtrise des outils d'IA est une compétence recherchée dans
la chasse aux bogues moderne. Utilisez-les sagement, citez-les
honnêtement, et conservez toujours votre propre compréhension de
chaque exploit que vous déployez.}


\chapter{Avant-propos}

nf\_tables est un terrain de chasse idéal pour les variantes parce que le
même objet logique est manipulé par plusieurs chemins :

\begin{itemize}
  \item le chemin de transaction, qui prépare l'état futur ;
  \item le chemin de commit, qui rend l'état visible ;
  \item le chemin d'abort, qui défait une transaction ;
  \item le chemin de GC, qui libère les objets après un délai de grâce.
\end{itemize}

Une variante naît quand un de ces chemins utilise une règle de validité qui
appartient à un autre. La question centrale est donc :

\begin{quote}
« Cette fonction vérifie-t-elle l'état courant ou l'état suivant, et
pourquoi ? »
\end{quote}

\chapter{Aperçu du module}

\begin{center}
\begin{tabularx}{\linewidth}{|c|l|c|c|}
\hline
\textbf{Ex} & \textbf{Titre} & \textbf{Difficulté} & \textbf{Points} \\ \hline
00 & Modèle de transaction et de génération & Facile & 3 \\
01 & Variantes de CVE-2023-35001 & Moyen & 5 \\
02 & Catchall, GC et double-free pipapo & Difficile & 5 \\
03 & Écrire les règles & Difficile & 5 \\
04 & Lancer et trier & Moyen & 4 \\ \hline
& \textbf{Total} & & \textbf{22} \\ \hline
\end{tabularx}
\end{center}

\startexercices

\chapter{Exercise 00}
\extitle{Modèle de transaction et de génération}
\turnindir{ex00}
\exnumber{00}
\exfiles{model.md}
\exauthorize{grep, un éditeur de texte}
\makeheaderfiles

Dans \texttt{model.md}, construisez un modèle exécutable mental de
nf\_tables :

\begin{enumerate}
  \item les objets d'une transaction ;
  \item les générations et le \texttt{genmask} ;
  \item la différence entre \texttt{nft\_is\_active()} et
        \texttt{nft\_is\_active\_next()} ;
  \item le moment où un objet devient visible ;
  \item le moment où un objet peut être libéré.
\end{enumerate}

Ajoutez un tableau « qui voit quoi » avec trois colonnes : transaction en
cours, commit, GC.

\hint{Un objet peut être actif dans la génération courante et inactif dans
la génération suivante. C'est cette asymétrie qui rend les tests de
génération difficiles à relire.}

\chapter{Exercise 01}
\extitle{Variantes de CVE-2023-35001}
\turnindir{ex01}
\exnumber{01}
\exfiles{nft-byteorder.md}
\exauthorize{git, grep, un éditeur de texte}
\makeheaderfiles

Le correctif de CVE-2023-35001 contraint un accès dans
\texttt{nft\_byteorder\_eval()}. Cherchez les autres opérations qui lisent
ou écrivent des registres avec une hypothèse de taille ou d'index.

\subsection*{Mission}

Dans \texttt{nft-byteorder.md} :

\begin{enumerate}
  \item listez au moins cinq opérations voisines ;
  \item pour chacune, indiquez le type de registre, la taille lue et la
        borne appliquée ;
  \item identifiez celles qui partagent l'hypothèse du correctif ;
  \item proposez un test de validation qui ne nécessite pas d'exploit ;
  \item donnez un statut et un score de triage à chaque candidate.
\end{enumerate}

\chapter{Exercise 02}
\extitle{Catchall, GC et double-free pipapo}
\turnindir{ex02}
\exnumber{02}
\exfiles{gc-catchall.md}
\exauthorize{git, grep, un éditeur de texte}
\makeheaderfiles

Le correctif de CVE-2024-1085 remplace :

\begin{console}
- if (!nft_is_active(net, ext))
+ if (!nft_is_active_next(net, ext))
      continue;
\end{console}

dans \texttt{nft\_setelem\_catchall\_deactivate()}.

\subsection*{Mission}

Dans \texttt{gc-catchall.md} :

\begin{enumerate}
  \item expliquez pourquoi le test doit porter sur la génération suivante
        dans ce chemin ;
  \item décrivez la conséquence d'un \texttt{kfree()} sur un élément encore
        actif ;
  \item reliez ce correctif à CVE-2023-4244 et à la suppression de la
        marque « busy » ;
  \item cherchez trois autres fonctions qui appellent
        \texttt{nft\_is\_active()} sans tester \texttt{\_next} ;
  \item pour chacune, dites si le test courant est correct ou suspect.
\end{enumerate}

\subsection*{Variante 2024 : CVE-2024-26809}

Le double-free pipapo de CVE-2024-26809 est un second cas de durée de vie.
Il ne faut pas le confondre avec CVE-2024-1085 : le premier touche la
destruction d'un set pipapo, le second la désactivation d'un élément
catchall.

Dans \texttt{gc-catchall.md}, ajoutez une section qui répond à :

\begin{enumerate}
  \item pourquoi \texttt{priv->dirty} et \texttt{priv->clone} peuvent
        provoquer une double libération ;
  \item pourquoi \texttt{set->dead} rend le chemin de destruction
        atteignable ;
  \item quel commit a introduit le code fautif, et quel commit a introduit
        le drapeau qui le rend exploitable ;
  \item pourquoi ce cas est une interaction entre deux hypothèses, et non
        une simple copie du correctif de CVE-2024-1085.
\end{enumerate}

\begin{warn}
Le chemin GC est souvent exécuté de manière asynchrone. Une lecture de code
linéaire peut donner l'illusion d'une exclusion mutuelle qui n'existe pas.
Ne reproduisez pas le double-free : l'exercice porte sur l'archéologie du
patch et sur la propriété des objets.
\end{warn}

\chapter{Exercise 03}
\extitle{Écrire les règles}
\turnindir{ex03}
\exnumber{03}
\exfiles{rules/}
\exauthorize{Semgrep, Weggli, un éditeur de texte}
\makeheaderfiles

Écrivez au moins trois règles :

\begin{enumerate}
  \item une règle Semgrep qui repère un appel à
        \texttt{nft\_is\_active()} dans une fonction dont le nom contient
        \texttt{deactivate} ou \texttt{destroy} ;
  \item une règle Weggli qui repère un \texttt{kfree()} proche d'un test
        d'activité ;
  \item une règle qui repère une fonction de désactivation sans test de
        génération.
\end{enumerate}

Chaque règle doit avoir :

\begin{itemize}
  \item un identifiant lisible ;
  \item un message qui explique \emph{pourquoi} le motif est suspect ;
  \item un exemple positif et un exemple négatif dans
        \texttt{rules/tests/}.
\end{itemize}

\chapter{Exercise 04}
\extitle{Lancer et trier}
\turnindir{ex04}
\exnumber{04}
\exfiles{findings.md}
\exauthorize{Semgrep, Weggli, un éditeur de texte}
\makeheaderfiles

Lancez vos règles sur \texttt{net/netfilter/} uniquement. Produisez
\texttt{findings.md} avec :

\begin{enumerate}
  \item le nombre de résultats bruts ;
  \item le nombre de résultats après déduplication ;
  \item cinq résultats triés, avec fichier, ligne et justification ;
  \item les faux positifs les plus instructifs ;
  \item une conclusion sur la précision de vos règles ;
  \item une comparaison avec le patch gap OtterRoot : commit public,
        publication de l'exploit KernelCTF, puis arrivée dans Ubuntu et
        Debian.
\end{enumerate}

Rappel de triage : score \texttt{>= 6} = investiguer, score
\texttt{>= 8} = prioriser. Un résultat sans invariant partagé est un faux
positif, même s'il ressemble au motif.

\chapter{Rendu et évaluation}

\begin{console}
A-EPI-002-nftables-variant-hunting/
|-- ex00/model.md
|-- ex01/nft-byteorder.md
|-- ex02/gc-catchall.md
|-- ex03/rules/
+-- ex04/findings.md
\end{console}

\begin{center}
\begin{tabularx}{\linewidth}{|l|X|c|}
\hline
\textbf{Critère} & \textbf{Attendu} & \textbf{Points} \\ \hline
Modèle & Génération et visibilité expliquées & 4 \\ \hline
Variantes & Cinq opérations comparées, statuts justifiés & 5 \\ \hline
GC & Lien catchall, busy mark et double free & 5 \\ \hline
Règles & Trois règles testées, exemples positifs/négatifs & 4 \\ \hline
Triage & Résultats dédupliqués et priorisés & 4 \\ \hline
\end{tabularx}
\end{center}

\chapter{Ressources par exercice}

\subsection*{Exercice 00 --- Modèle de transaction et de génération}
\begin{itemize}
  \item nf\_tables :
        \url{https://www.kernel.org/doc/html/latest/networking/nf_tables.html}
  \item Papier, sections 6.1 et 6.2.
  \item À chercher : \texttt{nft\_trans}, \texttt{genmask},
        \texttt{nft\_is\_active\_next}, \texttt{commit\_list}.
\end{itemize}

\subsection*{Exercice 01 --- Variantes de CVE-2023-35001}
\begin{itemize}
  \item NVD : \url{https://nvd.nist.gov/vuln/detail/CVE-2023-35001}
  \item Analyse : \url{https://pwning.tech/nftables/}
  \item Papier, section 6.2.
  \item À chercher : \texttt{nft\_byteorder}, \texttt{union},
        \texttt{priv->len}, \texttt{sizeof}, \texttt{stride}.
\end{itemize}

\subsection*{Exercice 02 --- Catchall, GC et double-free pipapo}
\begin{itemize}
  \item NVD : CVE-2023-4244, CVE-2024-1085, CVE-2024-26809.
  \item OtterSec, OtterRoot :
        \url{https://osec.io/blog/netfilter-universal-root-1-day/}
  \item Commits \texttt{b0e256f3dd2ba6532f37c5c22e07cb07a36031ee},
        \texttt{9827a0e6e23bf43003cd3d5b7fb11baf59a35e1e} et
        \texttt{5f68718b34a531a556f2f50300ead2862278da26}.
  \item À chercher : \texttt{nft\_set\_pipapo\_destroy},
        \texttt{nft\_set\_pipapo\_match\_destroy}, \texttt{priv->dirty},
        \texttt{set->dead}.
\end{itemize}

\subsection*{Exercice 03 --- Écrire les règles}
\begin{itemize}
  \item Semgrep : \url{https://semgrep.dev/docs/}
  \item Weggli : \url{https://github.com/weggli-rs/weggli}
  \item Dépôt compagnon : \texttt{rules/semgrep/} et \texttt{rules/cocci/}.
  \item À chercher : \texttt{nft\_is\_active}, \texttt{nft\_is\_active\_next},
        \texttt{scalar32\_min\_max}, \texttt{READ\_ONCE}.
\end{itemize}

\subsection*{Exercice 04 --- Lancer et trier}
\begin{itemize}
  \item Papier, section 18.8.
  \item Dépôt compagnon : \texttt{rules/TRIAGE.md}.
  \item NVD : \url{https://nvd.nist.gov/}
  \item À chercher : \texttt{false positive}, \texttt{twin},
        \texttt{reachable}, \texttt{score >= 6}.
\end{itemize}

\chapter{Sources}
\begin{itemize}
  \item Phrack, \emph{Security is failure: Every Patch Is a Guess},
        sections 6.2, 6.3 et 14.
  \item NVD : CVE-2023-35001, CVE-2023-4244, CVE-2024-1085,
        CVE-2024-26809.
  \item Commit \texttt{b1db244ffd041a49ecc9618e8feb6b5c1afcdaa7}.
  \item OtterSec, \emph{OtterRoot: Netfilter universal root 1-day} :
        \url{https://osec.io/blog/netfilter-universal-root-1-day/}
  \item Commits \texttt{b0e256f3dd2ba6532f37c5c22e07cb07a36031ee},
        \texttt{9827a0e6e23bf43003cd3d5b7fb11baf59a35e1e} et
        \texttt{5f68718b34a531a556f2f50300ead2862278da26}.
  \item Documentation nf\_tables et RCU.
\end{itemize}

\end{document}
